\cotefile{13}
\scanpage{11}
tout $s\in S$, $P_s$ le \unsure{torseur|torseurs} sur $G(\unsure{\mathbb{Q}|Q})$
\[
 P_s=\operatorname{Isom}_{\otimes}(\omega\inserted{\text{foncteur oubli}},\unsure{\Phi_s}F)
\]
Pour $s$ variable, les $P_s$ forment un revêtement principal \inserted{(à droite)} $P$ de $S$, de groupe $\Gamma=G(\mathbb{Q})$. On a tout ce qu'il fallait pour que le \unsure{$\otimes$} composé
\diagram{\begin{tikzcd}[ampersand replacement=\&,column sep=large]
 C=\operatorname{Rep}(G) \arrow[r,"F"] \& \operatorname{Hodge}(S) \arrow[r,"{V\mapsto V\times_S P}"] \& \operatorname{Hodge}(P)
\end{tikzcd}}
\deleted{\unsure{ait une struct|ait un struct}} \unsure{definisse|définisse} sur $P$ une $G$-\unsure{subordination}, d'où un morphisme
\[
 g:P\longrightarrow\unsure{\gamma|\mathcal{Y}}.
\]
\deleted{Ceci posé,} Ce dernier est compatible avec l'action de $\Gamma=G(\mathbb{Q})$\unsure{,|.} \inserted{(opérant sur $\gamma$ à \uline{droite}).} Et la catégorie des \deleted{\unsure{$\otimes$ foncteurs} $F$} $G$-subordinations $F$ sur $S$ compatibles avec le foncteur oubli est équivalente à celle des couples $(P,g)$, $P$ un rév. principal de $S$ de \unsure{gpe|groupe} $\Gamma$, $g:P\unsure{\relbar|\to}\gamma$ un $\Gamma$ morphisme.

Soit $\Gamma_0\subset\Gamma=G(\mathbb{Q})$ un \unsure{ss-gpe|ss-groupe} [de \unsure{préférence|préference|preférence}, commensurable à un $G(\mathbb{Z})$]. On appelle $G$-subordination $\Gamma_0$-\unsure{rigidifiée|rigidifiee} sur $S$ la donnée \unsure{d'une|d'un} rev. principal $P_0$ de $S$ de
% archive sheet number: 10
